\clearpage\begingroup
\SourceLeaf{320}\NoSourceHead\HistoricalFoot{}{40\textsuperscript{*}}
\begin{center}
{\large [ÜBER DEN\\REPSOLDSCHEN MERIDIANKREIS.]}\par\medskip
\textls[150]{BRIEFWECHSEL.}\par\medskip
\textsc{Gauss} an \textsc{Gerling}, Göttingen, 25. August 1818.
\end{center}

\noindent\hspace*{\parindent}. . . . . . Ueber den \textsc{Repsold}schen Kreis habe ich vor einigen Wochen einen kleinen Aufsatz in die hiesigen G[elehrten] A[nzeigen\begingroup\renewcommand{\thefootnote}{*)}\footnotemark\renewcommand{\thefootnote}{[*)}\footnotetext{1818 S. 1257, Werke VI, S. 410.]}\endgroup] gegeben, den Sie vermuthlich gelesen und daraus die Vortrefflichkeit dieses Instruments mit mehrerm ersehen haben werden. Die Mikroskope sind nicht unmittelbar an dem Stein, sondern an dem (schiebbaren Theile der) Lager befestigt, durch Arme, die sich auf zwei conischen Spitzen drehen, so dass die Mikroskope eine kleine Bewegung senkrecht auf die Fläche des Instruments haben. Dies ist deswegen nöthig, weil der Kreis keine absolut vollkommene Ebene ist, und mit Hülfe eines Röllchens an den Mikroskopen, welches auf der Kreisfläche läuft, behalten jene immer genau gleichen Abstand von dieser. Die Mikroskope werden nicht nach der Libelle eingestellt, sondern bleiben wo sie einmal sind; täglich mehreremal oder vielmehr nahe bei jeder wichtigen Beobachtung wird aber nivellirt und dabei die Mikroskope \textls[120]{abgelesen} (jedes 2mal, weil) allemal die Libelle auch umgehängt wird). Meine Erfahrung hat mich von der Nothwendigkeit dieses Verfahrens überzeugt; anfangs hatte ich seltner nivellirt und mehr auf die \textsc{Pond}sche Art observirt, indem ich die Sterne auf die nächsten Nordstern-Culminationen bezog; allein in 12 Stunden ändern sich die Mikroskope oft um 3 bis 4 Secunden und zwar immer gemeinschaftlich, welches auf ein Wanken des Steins, als Drehung um eine horizontale Axe von Ost nach West, zu deuten scheint. Ich halte es für einen Fehler des \textsc{Pond}schen Kreises, dass ihm ein solches Versicherungsmittel wie meine herrliche Libelle ist, fehlt.

Ein definitives Resultat für die Polhöhe werde ich erst nach Jahr und Tag geben können, wenn mehrere Circumpolarsterne in beiden Culminationen
\clearpage\SourceLeaf{321}\HistoricalHead{316}{L}{PRAKTISCHE UND SPHÄRISCHE ASTRONOMIE. REPSOLDSCHER MERIDIANKREIS.}
\noindent oft genug beobachtet sind. Vermuthlich liegt sie zwischen \(51^\circ31'49''\) und \(50''\). Ich theile Ihnen einige bis jetzt reducirte Zenithdistanzen auf den Anfang von 1818 mit, denen ich \textsc{Pond}s Polardistanzen und die daraus folgende Polhöhe beisetze. Natürlich ist dies Verfahren des \textsc{Repsold}schen Kreises unwürdig, und umgekehrt, wenn erst aus C[ircum]p[olar]sternen die Polhöhe abgeleitet ist, werden daraus die Declinationen geschlossen werden müssen. Indessen lässt sich vorläufig hieraus abnehmen, dass meine Declinationen von den \textsc{Pond}schen wohl nur sehr wenig abweichen werden. Bei \(\alpha\) Librae würde die Uebereinstimmung mit den übrigen noch grösser sein, wenn ich dieselbe Refraction wie \textsc{Pond} gebraucht hätte (ich habe \textsc{Bessel}s Tafel angewandt). Die Beob[achtungen] von mehrern andern Sternen habe ich noch nicht reducirt. Die Zahl der Beob[achtungen] ist nach der Formel \(\frac{4ab}{a+b}\) berechnet, wo \(a\) die Zahl der Beob[achtungen] bedeutet, wo der Kreis im Osten war, \(b\) die wo er im Westen war (gemäss meiner Wahrscheinlichkeitstheorie).

Lauter Tag-Beobachtungen
\begin{center}\begingroup\fontsize{9.5}{12}\selectfont\setlength{\tabcolsep}{2.4pt}
\begin{tabular}{@{}l|r|rrr|rrr|rrr@{}}\hline\hline
&Beobb.&\multicolumn{3}{c|}{Z[enith-]D[istanz]}&\multicolumn{3}{c|}{\shortstack{\textsc{Pond}s Pol[ar-]\\Dist[anz]}}&\multicolumn{3}{c}{Polhöhe}\\\hline\hline
Nordstern U. C.&\(16{,}5\)&\(-40^\circ\)&\(7'\)&\(53\overset{\prime\prime}{,}60\)&\(1^\circ\)&\(39'\)&\(44\overset{\prime\prime}{,}35\)&\(51^\circ\)&\(31'\)&\(50\overset{\prime\prime}{,}75\)\\
\(\beta\) Ursae min. O. C.&\(13{,}7\)&\(-23\)&22&\(7{,}18\)&15&6&\(2{,}50\)&&&\(50{,}32\)\\
Capella O. C.&\(8{,}9\)&\(+5\)&43&\(46{,}87\)&44&11&\(57{,}65\)&&&\(49{,}22\)\\
Arcturus&\(18{,}5\)&\(+31\)&23&\(43{,}25\)&69&51&\(53{,}95\)&&&\(49{,}30\)\\
\(\beta\) Leonis&14&\(+35\)&56&\(26{,}01\)&74&24&\(37{,}50\)&&&\(48{,}51\)\\
\(\alpha\) Orionis&\(8{,}9\)&\(+44\)&9&\(57{,}13\)&82&38&\(8{,}85\)&&&\(48{,}28\)\\
\(\alpha^2\) Librae&\(13{,}7\)&\(+66\)&48&\(30{,}11\)&105&16&\(38{,}70\)&&&\(51{,}41\)
\end{tabular}\par\endgroup\end{center}
\noindent . . . . . . . . . . . .

\begin{center}BEMERKUNGEN.\end{center}
\begingroup\small
\noindent\hspace*{\parindent}Die im Nachlass vorhandenen Stücke über \textsc{Gauss}’ Untersuchungen am \textsc{Repsold}schen Meridiankreise sind äusserst spärlich und zum Abdruck nicht geeignet. \textsc{Gauss} erhielt den Kreis im April 1818 und führte im wesentlichen neben der sorgfältigen Untersuchung des Instruments eine längere Reihe von Deklinationsbestimmungen von Zirkumpolarsternen und einige Rektascensionsbestimmungen, im besonderen des Polarsterns, aus. Die Beobachtungen mit dem Instrumente traten jedoch bald in den Hintergrund, als \textsc{Gauss} im August 1819 in den Besitz des \textsc{Reichenbach}schen Meridiankreises gelangte, dem er nun seine Aufmerksamkeit zuwandte. Vgl. den Aufsatz über \textsc{Gauss}’ astronomische Arbeiten, Werke XI2.
\begin{flushright}\textsc{Brendel}.\end{flushright}
\begin{center}\rule{.24\textwidth}{.3pt}\end{center}
\par\endgroup
\clearpage\endgroup
